%% titre: Méthode — Démontrer par récurrence
%% niveau: Licence 1
%% matiere: Raisonnement
%% annee: 2025-2026
%% ordre: 1
%% description: Exemple de fiche méthode (à remplacer) : les étapes d'une récurrence simple.

\section{Quand l'utiliser ?}

Pour démontrer une propriété $\mathcal{P}(n)$ valable pour tout entier $n\ge n_0$, notamment quand $\mathcal{P}(n+1)$ se déduit facilement de $\mathcal{P}(n)$ (sommes, suites définies par récurrence, divisibilité).

\section{La méthode}

\begin{methode}[Récurrence simple]
\begin{enumerate}[label=\textbf{Étape \arabic*.}]
  \item \textbf{Énoncer} clairement $\mathcal{P}(n)$, en précisant pour quels $n$ elle est visée.
  \item \textbf{Initialisation :} vérifier $\mathcal{P}(n_0)$.
  \item \textbf{Hérédité :} soit $n\ge n_0$ fixé tel que $\mathcal{P}(n)$ est vraie ; montrer $\mathcal{P}(n+1)$ en utilisant l'hypothèse de récurrence.
  \item \textbf{Conclure} : $\mathcal{P}(n)$ est vraie pour tout $n\ge n_0$.
\end{enumerate}
\end{methode}

\begin{attention}
Ne jamais écrire « supposons que $\mathcal{P}(n)$ est vraie pour tout $n$ » : on la suppose pour \emph{un} $n$ fixé.
\end{attention}

\section{Exemple traité}

\begin{exemple}
Montrer que, pour tout $n\in\N$, $\displaystyle\sum_{k=0}^{n}k=\frac{n(n+1)}{2}$.

\textbf{Initialisation.} Pour $n=0$ : $0=\frac{0\times 1}{2}$.

\textbf{Hérédité.} Si la formule est vraie au rang $n$, alors
\begin{align*}
\sum_{k=0}^{n+1}k &= \frac{n(n+1)}{2}+(n+1)\\
&= \frac{(n+1)(n+2)}{2},
\end{align*}
c'est la formule au rang $n+1$.

\textbf{Conclusion.} La formule est vraie pour tout $n\in\N$.
\end{exemple}

\begin{exercice}[À vous]
Montrer que $2^n\ge n+1$ pour tout $n\in\N$.
\begin{correction}
Pour $n=0$ : $1\ge 1$. Si $2^n\ge n+1$, alors $2^{n+1}\ge 2n+2\ge n+2$.
\end{correction}
\end{exercice}
